\subsection{Fine-Tuning Generative Robot Policies}

Generative policies require specialized reinforcement learning updates
because their action likelihoods are often intractable. DPPO optimizes a
diffusion policy through its denoising process~\cite{ren2024dppo}, while
ReinFlow introduces stochastic transitions into flow sampling to enable
policy-gradient updates~\cite{reinflow2025}. Flow Matching Policy Gradients
uses differences in flow-matching losses to construct a policy-update
surrogate~\cite{mcallister2025flowmatchingpolicygradients}. These methods
exploit the structure of the generator or its training objective. In
contrast, \method{} needs only samples from a frozen noise-to-action map
and a critic. This makes its residual update applicable to one-step
generators as well as iterative samplers. Our focus is how to preserve
pretrained behavior during critic-guided improvement.

\subsection{Critic-Guided Action Improvement}

Using a critic to improve actions before fitting a policy is an established
approach. DIPO modifies replay-buffer actions by gradient ascent and trains
a diffusion policy on the resulting targets~\cite{dipo2024}; Q-Score
Matching trains a diffusion policy using the critic's action
gradient~\cite{psenka2024qsm}. Particle-based methods provide another
connection: Soft Q-Learning trains an amortized sampler using Stein
variational updates~\cite{haarnoja2017sql}, and drifting models train
one-step generators by fitting displaced samples~\cite{deng2026drifting}.
Drifting Field Policy uses critic-ranked candidates to guide updates on a
drifting backbone~\cite{dfp2026}. \method{} builds on these ideas by
constructing fresh action targets for a residual on a frozen base, with
displacement clipping and explicit restoration toward the base. Stored
transitions are left unchanged.

An alternative is to use critic values without following their action
gradients. AWR and MPO favor actions through advantage- or value-weighted
fitting~\cite{peng2019awr,abdolmaleki2018mpo}, while IDQL combines a
diffusion behavior model with critic-based action
selection~\cite{hansenestruch2023idql}. These approaches motivate our
ranking-only and value-weighted field variants. We compare them with the
gradient field to distinguish the value of the improvement signal from the
effect of anchoring.

\subsection{Residual Learning and Pretrained-Policy Anchoring}

Residual learning preserves a base controller while learning corrections.
DICE-RL applies this construction to generative policies, combining a
critic-driven residual actor with behavior-cloning
regularization~\cite{dicerl2026}. It is our closest system-level reference
and the source of our controlled actor comparison. DAWN studies the critic
cold-start and scaling issues that affect residual RL~\cite{dawn2026}.
More broadly, TD3 and BCQ show how function-approximation and extrapolation
errors can undermine critic-guided optimization~\cite{fujimoto2018td3,fujimoto2019bcq}.

Our contribution is a controlled study of the restoring constraint: limiting
each displacement, pulling toward a fixed base, and allowing a region of
small corrections have different effects. The successful loss-space hinge
control is important to this positioning. It shows that useful restoring
geometry can be expressed within a conventional actor loss as well as through
\method{}'s action targets. We therefore study when that geometry preserves
and improves pretrained skills, rather than claiming that target regression
alone resolves critic error.
